\section{结论与展望}\label{sec:conclusion}

本文提出 AgriOcc，用三路前视图像直接预测农业机器人坐标系下的前方当前三维语义占用。该框架以预训练图像编码器、BEV 查询和占用解码器为基础，并通过 GVAD 以相机几何可见性构造全局锚点，通过 Agri-AMoE 将最终 BEV 表征提升为梯度能量路由的三维混合专家体积，通过近远场查询布局在保持近场稠密感知的同时优化早期编码计算。指定的 GVAD+Agri-AMoE+NearFar 运行在 FarmSim 上达到 56.99\% mIoU，较单帧 TSA 基线提升 1.14 个百分点；ORAD-3D 迁移结果进一步表明，FarmSim 预训练能够为异构真实越野类别空间提供有效的三维占用初始化。

AgriOcc 将农业环境中的作物、地面、自由空间和障碍物统一编码为可直接供规划使用的三维语义接口。其研究范式可自然扩展至更广泛的农业移动机器人任务：以多时刻观测增强动态作业场景建模，以选择性高分辨率子查询提升株间细缝与植株结构的表达，以语义占用为中间状态连接通行性预测、避障和闭环规划。

未来将进一步面向真实农田的光照、天气、作物生长期与传感器扰动构建更丰富的跨域训练策略；引入选择性激活的高分辨率近场查询，以提升株级作物间隙和细粒度边界的表达；并将 AgriOcc 的占用表征与风险感知通行性分析、轨迹优化和闭环导航深度耦合。最终，语义占用有望成为连接农业机器人感知、决策与控制的稳定三维基础模型接口。
